\begin{ouverture}[probabilites]{11}{Probabilités : comprendre le hasard}{Partie 3 · Données et hasard}
{Issues · Événements · Équiprobabilité · Contraire}
\objectif{Décrire une expérience aléatoire et ses issues.}
\objectif{Calculer une probabilité avec un modèle explicite.}
\objectif{Utiliser l’événement contraire et différentes écritures.}
\prerequis{Fractions, pourcentages et aires simples.}
\end{ouverture}
\pagestyle{otMain}\essentiel\fontsize{11}{15}\selectfont
\phantomsection\label{cours-11-1}
\essSection{1}{Décrire les possibilités}
\begin{definition}Une expérience aléatoire a plusieurs résultats possibles sans que l’on puisse prévoir avec certitude celui qui se produira. Chaque résultat possible est une issue. Un événement regroupe des issues qui répondent à une condition.\end{definition}
\begin{propriete}Une probabilité est un nombre compris entre 0 et 1. L’événement impossible a pour probabilité 0 ; l’événement certain a pour probabilité 1. La somme des probabilités de toutes les issues distinctes vaut 1.\end{propriete}
\begin{propriete}Si toutes les issues ont la même probabilité, on parle d’équiprobabilité. La probabilité d’un événement vaut alors le nombre d’issues favorables divisé par le nombre total d’issues. Cette formule nécessite de justifier pourquoi les issues comptées ont la même chance.\end{propriete}
\begin{exemple}Dans un sac de 3 boules rouges et 2 bleues, indiscernables au toucher et tirées au hasard, la probabilité de rouge vaut $3/5$. Les deux couleurs ne sont pas équiprobables ; les cinq boules le sont.\end{exemple}
\clearpage
\phantomsection\label{cours-11-2}
\essSection{2}{Utiliser un modèle cohérent}
\begin{definition}L’événement contraire de A regroupe toutes les issues qui ne réalisent pas A. Il se note ici « non A ». Les deux probabilités ont pour somme 1 : $P(\text{non A})=1-P(A)$.\end{definition}
\begin{exemple}Si la probabilité de gagner vaut 0,35 et s’il n’y a que gain ou perte, la probabilité de perdre vaut $1-0{,}35=0{,}65$, soit 65\,\%. Pour un dé équilibré, « ne pas obtenir 6 » a pour probabilité $5/6$.\end{exemple}
\label{fig-cours-11-2-cours-urne}\figureGeo{cours-urne}

\begin{savoirfaire}[Cas non équiprobable]Les probabilités peuvent être données par un modèle, par exemple 0,2 pour une face et 0,8 pour une autre. Pour une cible où le point d’impact est uniforme sur toute l’aire, une probabilité se calcule comme un rapport d’aires.\end{savoirfaire}
\begin{attention}Le résultat d’un essai ne change pas une probabilité annoncée pour des essais indépendants dans les mêmes conditions. Une fréquence observée est un résultat expérimental, pas une garantie sur le prochain tirage. Les calculs de ce chapitre portent sur une seule épreuve.\end{attention}

\clearpage\rubrique{Découverte}
\begin{blocLivre}{ex-11-1}\begin{exercice}[1]{Un dé équilibré}
On lance un dé équilibré numéroté de 1 à 6. Énumérer les issues, puis calculer la probabilité d’obtenir 4.
\end{exercice}\end{blocLivre}
\begin{blocLivre}{ex-11-2}\begin{exercice}[1]{Un événement}
Sur ce dé, calculer la probabilité d’obtenir un nombre strictement supérieur à 4.
\end{exercice}\end{blocLivre}
\begin{blocLivre}{ex-11-3}\begin{exercice}[1]{Certain et impossible}
Sur un dé classique, donner un événement certain et un événement impossible.
\end{exercice}\end{blocLivre}
\begin{blocLivre}{ex-11-4}\begin{exercice}[1]{Plusieurs écritures}
Une probabilité vaut $3/8$. L’écrire sous forme décimale et en pourcentage.
\end{exercice}\end{blocLivre}
\begin{blocLivre}{ex-11-5}\begin{exercice}[1]{Le contraire}
Un jeu a une probabilité de gain de 0,28. Il ne peut y avoir que gain ou perte. Quelle est la probabilité de perte ?
\end{exercice}\end{blocLivre}

\rubrique{Entraînement}
\begin{blocLivre}{ex-11-6}\begin{exercice}[2]{Une urne}
Une urne contient 4 rouges, 3 bleues et 1 verte, indiscernables au toucher. On tire une boule au hasard. Calculer les probabilités de rouge et de non-rouge.
\label{fig-ex-11-6}\figureGeo{urne-431}
\end{exercice}\end{blocLivre}
\begin{blocLivre}{ex-11-7}\begin{exercice}[2]{Une roue}
Une roue comporte 8 secteurs de même angle, numérotés de 1 à 8 et supposés équiprobables. Calculer la probabilité d’obtenir un multiple de 3.
\label{fig-ex-11-7}\figureGeo{roue-8}
\end{exercice}\end{blocLivre}
\begin{blocLivre}{ex-11-8}\begin{exercice}[2]{Des faces inégales}
Un objet a trois issues A, B, C de probabilités respectives 0,2 ; 0,5 ; $p$. Trouver $p$ et la probabilité de ne pas obtenir B.
\end{exercice}\end{blocLivre}
\begin{blocLivre}{ex-11-9}\begin{exercice}[2]{Une cible carrée}
Un point est choisi uniformément sur une cible carrée de côté 10 cm. Une zone gagnante est un carré de côté 4 cm entièrement à l’intérieur. Calculer la probabilité de gagner.
\label{fig-ex-11-9}\figureGeo{cible-carree}
\end{exercice}\end{blocLivre}
\begin{blocLivre}{ex-11-10}\begin{exercice}[2]{Une observation}
Sur 50 lancers d’une pièce annoncée équilibrée, on observe 31 piles. Quelle est la fréquence ? Doit-on annoncer que la probabilité du prochain pile vaut 31/50 ?
\end{exercice}\end{blocLivre}

\rubrique{Problèmes}
\begin{blocLivre}{ex-11-11}\begin{exercice}[3]{Un effectif inconnu}
Un sac contient 5 boules rouges et un nombre $b$ de bleues. Un tirage au hasard donne une probabilité de rouge égale à $1/4$. Trouver $b$.
\end{exercice}\end{blocLivre}
\begin{blocLivre}{ex-11-12}\begin{exercice}[3]{Une couronne gagnante}
Une cible est un disque de rayon 6 cm. La zone gagnante est la couronne extérieure au disque concentrique de rayon 3 cm. Le point d’impact est uniforme. Calculer la probabilité de gagner.
\label{fig-ex-11-12}\figureGeo{cible-couronne}
\end{exercice}\end{blocLivre}
\begin{blocLivre}{ex-11-13}\begin{exercice}[3]{Changer la composition}
Une urne contient 3 rouges et 5 bleues. On ajoute seulement des rouges pour que la probabilité de rouge devienne $1/2$. Combien faut-il en ajouter ?
\end{exercice}\end{blocLivre}
\begin{blocLivre}{ex-11-14}\begin{exercice}[3]{Une roue non uniforme}
Une roue circulaire a trois secteurs de $90^\circ$, $120^\circ$ et $150^\circ$. L’arrêt est uniforme en angle. Calculer la probabilité du dernier secteur, puis celle de ne pas l’obtenir.
\label{fig-ex-11-14}\figureGeo{roue-inegale}
\end{exercice}\end{blocLivre}
\begin{blocLivre}{ex-11-15}\begin{exercice}[3]{Des numéros et des couleurs}
Un sac contient six jetons équiprobables : R1, R2, R3, B1, B2, V1. Calculer la probabilité de « rouge ou portant le numéro 1 ». Pourquoi ne peut-on pas additionner directement les deux probabilités ?
\label{fig-ex-11-15}\figureGeo{jetons-numerotes}
\end{exercice}\end{blocLivre}

\clearpage
